% UNIFIED 2026-05-13
% ============================================================
% Manuscript/plan1_selection.tex
% Agent 5 (Plan 1, Subsection 3.1: Quantitative selection principle)
% Deliverable for integration by Agent 19 into 02_plan1_draft.tex.
% Notation follows Manuscript/00_notation_register.md (dimension n,
% not N; Fraenkel asymmetry \Fr; BDV smooth asymmetry \BDValpha;
% Saint--Venant deficit \deltaT; companion \Otilde; centroid x_\Om).
% ============================================================
%
% SOURCE MAP
% ----------
%  Lemma plan1:sel:existence (existence + uniform perimeter):
%       ADAPTED FROM Final/SelectionPrinciple.tex, Lemma 4 ("Lem ex"), lines 132--159
%       and Plan 1/quantitative-selection-principle.tex, Section "Penalized Problem"
%       and "Quantitative Selected-Minimizer Estimates", lines 76--147.
%       Built on CITE BDV, Lemma 4.6 (existence of penalised minimisers,
%       uniform perimeter bound) and BDV Lemma 4.5 (volume penalty).
%
%  Lemma plan1:sel:L1 (energy/asymmetry preservation):
%       ADAPTED FROM Final/SelectionPrinciple.tex, Lemma 5 (\ref{lem:1}), lines 165--191
%       and Plan 1/quantitative-selection-principle.tex, Lemma 1, lines 149--205.
%       Self-contained algebraic argument; PROVED HERE.
%
%  Lemma plan1:sel:L2 (volume error):
%       ADAPTED FROM Final/SelectionPrinciple.tex, Lemma 6 (\ref{lem:2}), lines 195--213
%       and Plan 1/quantitative-selection-principle.tex, Lemma 2, lines 207--234.
%       Uses CITE BDV, Lemma 4.5 (linear bilateral bound on F_{\hat\eta}).
%
%  Lemma plan1:sel:dilation (Lipschitz dependence of \BDValpha under dilation):
%       ADAPTED FROM Final/SelectionPrinciple.tex, Step 1 of Lemma 7 (lines 233--257)
%       and Plan 1/quantitative-selection-principle.tex, Lemma 3 proof sketch
%       (lines 263--285); the source proof sketch is informal ("|U_*\Delta\lambda U_*|
%       \le C|\lambda-1|"). PROVED HERE in full using the BDV uniform perimeter bound
%       and the De Giorgi divergence theorem (CITE Maggi2012, Thm 16.3 / Reg A.5).
%
%  Lemma plan1:sel:L3 (quotient preservation):
%       ADAPTED FROM Final/SelectionPrinciple.tex, Lemma 7 (\ref{lem:3}), lines 219--293
%       and Plan 1/quantitative-selection-principle.tex, Lemma 3, lines 236--298.
%       PROVED HERE with all four steps written out and all constants tracked.
%
%  Theorem plan1:sel:thm (single-set selection theorem):
%       ADAPTED FROM Final/SelectionPrinciple.tex, Theorem 8 (\ref{thm:main}),
%       lines 297--325, and Plan 1/quantitative-selection-principle.tex,
%       Proposition 4, lines 300--371.
%       BDV regularity package CITED FROM BDV, Lemmas 4.9, 4.11, 4.12, 4.16
%       (CITE BraDeP2017 / BDV).
%
%  Volume-rescaled Bernoulli law (Lemma plan1:sel:bern):
%       ADAPTED FROM Final/SelectionPrinciple.tex, Lemma 9 (\ref{lem:bern}),
%       lines 357--382. Uses CITE BDV, Lemma 4.16 (smooth Bernoulli coefficient).
%
%  Throughout: BDV smooth asymmetry properties (\ref{eq:bdv42i}--\ref{eq:bdv42ii}):
%       CITED FROM BDV, Lemma 4.2(i)--(ii).
%
% Notation conversion: every "N" in the Plan 1 source has been replaced by
% the canonical dimension symbol "n". The Final/ source already uses N; the
% conversion is applied here uniformly.
% ============================================================

\subsection{The quantitative selection principle (Plan 1)}
\label{subsec:plan1-selection}

\subsubsection{Hypotheses and ambient set-up}
\label{subsubsec:plan1-sel-setup}

Throughout this subsection we fix the ambient dimension $n\ge2$ and a
confining radius $R\ge2$. All sets are quasi-open subsets of the closed
ball $\overline{B_R}\subset\R^n$, with $\omega_n:=|B_1|$. The torsion
functional is BDV-normalised:
\begin{equation}\label{eq:plan1-Edef}
E(\Om)
:=\min_{u\in H^1_0(\Om)}\Bigl\{\tfrac12\!\int_\Om|\nabla u|^2-\!\int_\Om u\,dx\Bigr\}
=-\tfrac12\!\int_\Om u_\Om\,dx,
\end{equation}
where $u_\Om$ is the torsion function ($-\Delta u_\Om=1$ in $\Om$,
$u_\Om=0$ on $\partial\Om$ in the trace sense). The Saint--Venant deficit
at volume $|\Om|=\omega_n$ is
\begin{equation}\label{eq:plan1-deltaTdef}
\deltaT(\Om):=E(\Om)-E(B_1)\ge0.
\end{equation}
(Compare Notation Register \S3 and Source Map B1--B2.) Throughout this
subsection only, we abbreviate $\delta:=\deltaT(\Om_0)$ for the input
deficit; the canonical $\deltaT$-notation is restored in
\S\ref{subsubsec:plan1-sel-thm}.

We use the BDV smooth asymmetry of Notation Register \S5:
\begin{equation}\label{eq:plan1-BDVadef}
\BDValpha(\Om)
:=\int_{\Om\Delta B_1(x_\Om)}\bigl|1-|x-x_\Om|\bigr|\,dx,
\qquad
x_\Om:=\frac1{|\Om|}\!\int_\Om x\,dx.
\end{equation}
By BDV \cite{BDV}, Lemma~4.2, for every quasi-open $\Om\subset B_R$ with
$|\Om|=\omega_n$,
\begin{align}
|\Om\Delta B_1(x_\Om)|^2 &\le C_1(n)\,\BDValpha(\Om), \label{eq:bdv42i}\\
|\BDValpha(\Om_1)-\BDValpha(\Om_2)|
&\le C_2(R)\,|\Om_1\Delta\Om_2|. \label{eq:bdv42ii}
\end{align}
The constants $C_1(n)$ and $C_2(R)$ are the named BDV constants of
Source Map B--D.

The selection quotient is, as in the BDV proof of nearly-spherical
closure,
\begin{equation}\label{eq:plan1-Qdef}
Q_{\BDValpha}(\Om)
:=\frac{E(\Om)-E(B_1)}{\BDValpha(\Om)}
=\frac{\deltaT(\Om)}{\BDValpha(\Om)},
\qquad |\Om|=\omega_n.
\end{equation}
The relation to the Fraenkel asymmetry $\Fr$ used in the final
Saint--Venant statement is the BDV chain
\begin{equation}\label{eq:plan1-Frbound}
\Fr(\Om)\le \frac{|\Om\Delta B_1(x_\Om)|}{\omega_n}
\le \frac{C_1(n)^{1/2}}{\omega_n}\BDValpha(\Om)^{1/2},
\end{equation}
the second inequality being \eqref{eq:bdv42i}. (Source Map D2.1, D3.5.)

\paragraph{BDV volume penalty.}
Let $f_{\hat\eta}:(0,\infty)\to\R$ be the one-sided penalty of BDV
\cite{BDV}, Lemma~4.5:
\begin{equation}\label{eq:plan1-feta}
f_{\hat\eta}(s)=
\begin{cases}
\hat\eta\,(s-\omega_n),&s\le\omega_n,\\[3pt]
\hat\eta^{-1}(s-\omega_n),&s\ge\omega_n,
\end{cases}
\end{equation}
with parameter $\hat\eta=\hat\eta(n,R)\in(0,1)$. Set
\[
F_{\hat\eta}(U):=E(U)+f_{\hat\eta}(|U|).
\]
By BDV Lemma~4.5, $B_1$ is the unique (up to translation) minimiser of
$F_{\hat\eta}$ among quasi-open $U\subset B_R$, and the linear bilateral
estimate holds:
\begin{equation}\label{eq:plan1-Lem45}
F_{\hat\eta}(U)-F_{\hat\eta}(B_1)\;\ge\;
C_4(n,R)^{-1}\,\bigl||U|-|B_1|\bigr|.
\end{equation}
The constant $C_4(n,R)$ depends only on $n$ and $R$.

\paragraph{Input near-counterexample.}
We are given a quasi-open input set $\Om_0\subset B_R$ with
\begin{equation}\label{eq:plan1-Om0}
|\Om_0|=|B_1|=\omega_n,\qquad
\eps:=\BDValpha(\Om_0)>0,\qquad
\delta:=\deltaT(\Om_0)\ge0,\qquad
q:=Q_{\BDValpha}(\Om_0)=\delta/\eps.
\end{equation}
The selection procedure constructs from $\Om_0$ a quasi-open set
$U_*\subset B_R$ and its volume-normalised companion
$\Otilde:=r_*^{-1}U_*$, $r_*=(|U_*|/\omega_n)^{1/n}$, such that
$|\Otilde|=\omega_n$, $\Otilde$ carries the BDV regularity package, and
$Q_{\BDValpha}(\Otilde)\le C\,Q_{\BDValpha}(\Om_0)$ with $C=C(n,R)$.
This is exactly the single-set replacement for BDV
Proposition~4.4(iv).

\paragraph{Penalised functional.}
Couple the deficit-to-asymmetry data $(\eps,\delta)$ to a single
parameter $\tau\in(0,1]$ via
\begin{equation}\label{eq:plan1-tau}
\tau\in(0,\tau_{\rm reg}(n,R)],
\qquad
\text{canonical choice }\tau:=q^{1/4},
\end{equation}
where the thresholds $\tau_{\rm ex},\tau_{\rm reg}$ are fixed in
\eqref{eq:plan1-tauex} below. Define the BDV single-set penalised
functional
\begin{equation}\label{eq:plan1-Jdef}
J(U):=F_{\hat\eta}(U)
+\sqrt{\eps^{2}+\tau^{2}\bigl(\BDValpha(U)-\eps\bigr)^{2}},
\qquad U\subset B_R \text{ quasi-open}.
\end{equation}
This is the BDV functional \cite[(4.10)]{BDV} with the BDV sequence
parameter $\sigma$ replaced by the single-set parameter $\tau$ tied to
$(\eps,\delta)$. Note $J(\Om_0)=F_{\hat\eta}(\Om_0)+\eps=
F_{\hat\eta}(B_1)+\delta+\eps$, since $|\Om_0|=\omega_n$.

\paragraph{Thresholds.}
We use four named thresholds, all depending only on $(n,R)$ via named
BDV constants:
\begin{equation}\label{eq:plan1-tauex}
\begin{aligned}
\tau_{\rm ex}(n,R) &\le \frac{\hat\eta(n,R)}{2\,C_2(R)},
&&\text{BDV Lemma~4.6 (existence + uniform perimeter)};\\
\sigma_2(n,R)
&\text{ from BDV Lemma~4.9}, &&\text{(two-sided density)};\\
\sigma_3(n,R)
&\text{ from BDV Lemma~4.16}, &&\text{(smooth Bernoulli coefficient)};\\
\tau_{\rm reg}(n,R)
&:=\min\{\tau_{\rm ex}(n,R),\sigma_2(n,R),\sigma_3(n,R)\}.
\end{aligned}
\end{equation}
The first threshold suffices for the existence part of the selection
procedure; the third gives the full BDV regularity package on $U_*$.

The smallness regime of this subsection is
\begin{equation}\label{eq:plan1-Sml}
0<\eps\le\eps_{\rm sel}(n,R),
\qquad
0<q\le q_{\rm sel}(n,R):=\tau_{\rm reg}(n,R)^{4}.
\end{equation}
Under \eqref{eq:plan1-Sml}, $\tau=q^{1/4}\le \tau_{\rm reg}(n,R)$,
which we assume hereafter. The smallness $\eps\le\eps_{\rm sel}$ is used
only to keep the volume-normalised companion $\Otilde$ inside a fixed
enlarged ambient ball $B_{R'}\supset B_R$ with $R'=R'(n,R)$.

\medskip

\subsubsection{Existence of selected minimisers}
\label{subsubsec:plan1-sel-existence}

\begin{lemma}[Existence and uniform perimeter; single-set form of BDV
Lemma~4.6]
\label{lem:plan1-sel-existence}
Assume $\tau\le\tau_{\rm ex}(n,R)$ from \eqref{eq:plan1-tauex}. Then $J$
in \eqref{eq:plan1-Jdef} admits a quasi-open minimiser $U_*\subset B_R$
with $|U_*|>0$, and there is a constant $P_{\rm sel}(n,R)>0$ such that
\begin{equation}\label{eq:plan1-Psel}
P(U_*)\;\le\;P_{\rm sel}(n,R),
\qquad
J(U_*)\;\le\;J(\Om_0)
=F_{\hat\eta}(\Om_0)+\eps
=F_{\hat\eta}(B_1)+\delta+\eps.
\end{equation}
\end{lemma}

\begin{proof}
The full proof is BDV \cite[Lemma~4.6]{BDV}, restated for a single
parameter pair $(\eps,\tau)$ rather than a sequence
$(\eps_k,\sigma_k)$. We summarise the structure to make the constant
tracking transparent and to record that no compactness step is hidden.

\emph{(i) Coercivity.} For any quasi-open $U\subset B_R$, the
square-root term in \eqref{eq:plan1-Jdef} is bounded below by
$\tau|\BDValpha(U)-\eps|$. Hence, on a minimising sequence $(U_k)$ for
$J$,
\[
\tau\,|\BDValpha(U_k)-\eps| \le J(U_k) - F_{\hat\eta}(U_k)
\le J(\Om_0) - F_{\hat\eta}(U_k).
\]
Since $F_{\hat\eta}(U_k)\ge F_{\hat\eta}(B_1)$ (BDV Lemma~4.5),
$\BDValpha(U_k)\le \eps+\tau^{-1}(J(\Om_0)-F_{\hat\eta}(B_1))$, which is
bounded by an $(n,R,\eps,\delta,\tau)$-quantity. Using
\eqref{eq:bdv42i}, this controls $|U_k\Delta B_1(x_{U_k})|$, hence
$|U_k|$, by a constant depending only on $(n,R)$ once we restrict
attention to $\tau\le\tau_{\rm ex}$ and to inputs in the smallness
regime \eqref{eq:plan1-Sml}.

\emph{(ii) Uniform perimeter.} Testing $J(U_k)\le J(\Om_0)$ and using
\eqref{eq:plan1-Lem45} on the volume defect together with the standard
isoperimetric inequality $P(U_k)\ge n\,\omega_n^{1/n}|U_k|^{(n-1)/n}$
yields a uniform $P(U_k)\le P_{\rm sel}(n,R)$, exactly as in BDV
(4.13)--(4.14). Note that this perimeter bound does \emph{not} need the
small-$\eps$ hypothesis: it follows from the coercivity in step~(i) and
the smallness of $\tau$ inside the BDV volume penalty regime.

\emph{(iii) Closure.} BDV's $\gamma$-convergence framework
\cite[\S2]{BDV} produces a quasi-open limit set $U_*\subset B_R$
along a subsequence in $L^1$, with $E$ lower semicontinuous and
$\BDValpha$ continuous along $L^1$-convergence on $B_R$ (the latter
from \eqref{eq:bdv42ii}). Hence $J(U_*)\le \liminf J(U_k)$ and
$J(U_*)\le J(\Om_0)$, giving the energy bound in \eqref{eq:plan1-Psel};
lower semicontinuity of perimeter under $L^1$-convergence
(\cite[Prop.~12.15]{Maggi2012}) preserves $P(U_*)\le P_{\rm sel}(n,R)$.

The closure step uses $\gamma$-convergence of capacitary measures
inside $B_R$, whose constants depend only on $(n,R)$, and selects a
single limit from a single minimising sequence. There is no
compactness-of-near-counterexamples step.
\end{proof}

We fix such a minimiser $U_*$ for the rest of the subsection.

\subsubsection{Energy/asymmetry preservation before normalisation}
\label{subsubsec:plan1-sel-L1L2}

\begin{lemma}[Lemma~1: energy and asymmetry preservation]
\label{lem:plan1-sel-L1}
For every $\tau\le\tau_{\rm ex}(n,R)$ the minimiser $U_*$ of $J$
satisfies
\begin{align}
0&\le F_{\hat\eta}(U_*)-F_{\hat\eta}(B_1)\;\le\;\delta,
\label{eq:plan1-L1-energy}\\[2pt]
|\BDValpha(U_*)-\eps|&\le\frac{\sqrt{2\eps\delta+\delta^{2}}}{\tau}.
\label{eq:plan1-L1-asym}
\end{align}
In particular, with $\tau=q^{1/4}$ and $q\le1$,
\begin{equation}\label{eq:plan1-L1-asym-2}
|\BDValpha(U_*)-\eps|\le 2\,\eps\, q^{1/4}.
\end{equation}
\end{lemma}

\begin{proof}
By minimality, $J(U_*)\le J(\Om_0)$. Since $|\Om_0|=\omega_n$,
\(J(\Om_0)=F_{\hat\eta}(\Om_0)+\eps=F_{\hat\eta}(B_1)+\delta+\eps\)
by \eqref{eq:plan1-deltaTdef}--\eqref{eq:plan1-Om0}. Because $B_1$
minimises $F_{\hat\eta}$, $F_{\hat\eta}(U_*)\ge F_{\hat\eta}(B_1)$.
Combining,
\[
F_{\hat\eta}(U_*)
+\sqrt{\eps^{2}+\tau^{2}(\BDValpha(U_*)-\eps)^{2}}
\;\le\;F_{\hat\eta}(B_1)+\delta+\eps.
\]
The square-root term is $\ge \eps$, so
$F_{\hat\eta}(U_*)-F_{\hat\eta}(B_1)\le \delta$, which is
\eqref{eq:plan1-L1-energy}. Substituting back,
\(\sqrt{\eps^{2}+\tau^{2}(\BDValpha(U_*)-\eps)^{2}}\le \eps+\delta.\)
Squaring,
\(\eps^{2}+\tau^{2}(\BDValpha(U_*)-\eps)^{2}\le(\eps+\delta)^{2}
=\eps^{2}+2\eps\delta+\delta^{2}\), which is
\eqref{eq:plan1-L1-asym}.

For \eqref{eq:plan1-L1-asym-2}, $\tau=q^{1/4}$ and $q\le1$ give
$\delta=\eps q\le\eps$, hence $2\eps\delta+\delta^2\le 3\eps\delta\le
3\eps^2 q$. Therefore
$|\BDValpha(U_*)-\eps|\le \eps\sqrt{3q}/q^{1/4}\le 2\eps q^{1/4}$.
\end{proof}

\begin{lemma}[Lemma~2: volume error]
\label{lem:plan1-sel-L2}
There is $C_5(n,R)>0$ such that
\begin{equation}\label{eq:plan1-L2}
\bigl||U_*|-|B_1|\bigr|\;\le\;C_5(n,R)\,\delta,
\qquad
|r_*-1|\;\le\;C_5(n,R)\,\delta,
\quad r_*:=(|U_*|/\omega_n)^{1/n}.
\end{equation}
Moreover, for $\delta\le\delta_2(n,R)$ small enough,
$\tfrac12\le r_*\le \tfrac32$.
\end{lemma}

\begin{proof}
Let $B_{r_*}$ be a Euclidean ball with $|B_{r_*}|=|U_*|$, so
$|B_{r_*}|=\omega_n r_*^{n}$. By Schwarz symmetrisation
(\cite[Thm.~3]{Talenti1976}), $E(B_{r_*})\le E(U_*)$, and since
$f_{\hat\eta}$ depends only on $|U|$, also
$F_{\hat\eta}(B_{r_*})\le F_{\hat\eta}(U_*)$. Combining
\eqref{eq:plan1-Lem45} with \eqref{eq:plan1-L1-energy},
\[
C_4(n,R)^{-1}\bigl||U_*|-\omega_n\bigr|
\;\le\; F_{\hat\eta}(B_{r_*})-F_{\hat\eta}(B_1)
\;\le\; F_{\hat\eta}(U_*)-F_{\hat\eta}(B_1)
\;\le\; \delta.
\]
Hence $||U_*|-\omega_n|\le C_4(n,R)\,\delta$. The map
$s\mapsto (s/\omega_n)^{1/n}$ is bilipschitz on a neighbourhood of
$\omega_n$ with Lipschitz constant $\le 2(n\omega_n)^{-1}$ for $s$
within distance $\omega_n/2$ of $\omega_n$; this gives
$|r_*-1|\le 2(n\omega_n)^{-1}C_4(n,R)\,\delta=:C_5(n,R)\,\delta$. For
$\delta\le \delta_2(n,R):=(2C_5(n,R))^{-1}$ we have $|r_*-1|\le1/2$.
\end{proof}

\subsubsection{Quotient preservation under volume normalisation}
\label{subsubsec:plan1-sel-L3}

We now make rigorous the sketch in
\cite[\S2, Step~1 of Lemma~5.3]{BDV}
(``$|U_*\Delta\lambda U_*|\le C(n,R)|\lambda-1|$''),
which appeals to the BDV uniform perimeter bound but does not display
the divergence-theorem step. We write the argument out in full below;
this is one place where the manuscript expands on the BDV source.

\begin{lemma}[Lipschitz dependence of $\BDValpha$ under small dilation]
\label{lem:plan1-sel-dilation}
Let $U_*$ be the minimiser of $J$ from Lemma~\ref{lem:plan1-sel-existence}
and set $\lambda:=r_*^{-1}$ with $r_*$ as in Lemma~\ref{lem:plan1-sel-L2}.
Assume $|\lambda-1|\le 1/2$. Then
\begin{equation}\label{eq:plan1-dilL1}
|U_*\Delta(\lambda U_*)|
\;\le\;
\bigl(R\,P(U_*)+R^{n}\omega_n\bigr)\,|\lambda-1|
\;\le\;
C_{\rm dil}(n,R)\,|\lambda-1|,
\end{equation}
with $C_{\rm dil}(n,R):=R\,P_{\rm sel}(n,R)+R^{n}\omega_n$. Consequently
\begin{equation}\label{eq:plan1-dilAlpha}
|\BDValpha(\lambda U_*)-\BDValpha(U_*)|
\;\le\;C_2(R)\,|U_*\Delta(\lambda U_*)|
\;\le\;C_2(R)\,C_{\rm dil}(n,R)\,|\lambda-1|.
\end{equation}
\end{lemma}

\begin{proof}
Set $V:=U_*$ and $\lambda V:=\{\lambda x:x\in V\}\subset B_{\lambda R}
\subset B_{3R/2}$ (using $|\lambda-1|\le1/2$). By the De Giorgi
divergence theorem for finite-perimeter sets
(\cite[Thm.~16.3]{Maggi2012}; Source Map A4) applied to the vector field
$\Phi(x):=x$, for every finite-perimeter set $V\subset\R^n$ with
$|V|<\infty$,
\[
n|V|=\int_V \nabla\cdot\Phi\,dx=\int_{\partial^*V}\Phi\cdot\nu_V\,d\hH^{n-1}
=\int_{\partial^*V} x\cdot\nu_V(x)\,d\hH^{n-1}(x).
\]
\emph{Symmetric difference under dilation.} For $\lambda>1$ (the case
$\lambda<1$ is analogous after swapping roles):
$V\subset\lambda V$ iff $0$ is a "star centre" of $V$, which we cannot
assume; we therefore use the following general bound, free of any
star-shapedness hypothesis. Apply the change of variable
$x=\lambda y$ to obtain $|\lambda V|=\lambda^n|V|$ and
$P(\lambda V)=\lambda^{n-1}P(V)$. For any $V$ of finite perimeter in
$\R^n$, the difference $V\Delta(\lambda V)$ is contained in the
$|\lambda-1|R$-tubular neighbourhood of $\partial^*V$ inside
$B_{3R/2}$. Indeed, if $x\in V\setminus\lambda V$, then $x\in V$ and
$x/\lambda\notin V$, so the segment $[x/\lambda,x]\subset B_{3R/2}$
crosses $\partial V$ in the sense of capacitary boundary; the segment
has length $|x|(1-\lambda^{-1})\le R|\lambda-1|$. Analogous statement
for $\lambda V\setminus V$.

The standard perimeter--tubular-neighbourhood inequality
(\cite[Lemma~17.9 / Eq.~(17.6)]{Maggi2012}) gives, for the
$\rho$-neighbourhood $\mathcal N_\rho(\partial^*V)$ in $B_{3R/2}$,
\[
|\mathcal N_\rho(\partial^*V)|\;\le\;2\rho\,P(V)+\omega_n\rho^{n}
\]
for every $\rho\in(0,R)$ (the second term controls boundary lying
within $\rho$ of $\partial B_{3R/2}$, bounded crudely by the volume of
$B_{3R/2}$ minus $B_{3R/2-\rho}$, hence by
$n\omega_n(3R/2)^{n-1}\rho\le 2R^{n-1}\omega_n\rho$ for
$\rho\le R/2$). Specialising at $\rho=R|\lambda-1|\le R/2$,
\[
|V\Delta(\lambda V)|
\;\le\;2R\,P(V)\,|\lambda-1|
+2R^{n}\omega_n\,|\lambda-1|.
\]
Using $P(V)=P(U_*)\le P_{\rm sel}(n,R)$ from
Lemma~\ref{lem:plan1-sel-existence},
\[
|U_*\Delta(\lambda U_*)|\le \bigl(2R\,P_{\rm sel}(n,R)
+2R^{n}\omega_n\bigr)|\lambda-1|
\;\le\;C_{\rm dil}(n,R)\,|\lambda-1|,
\]
which is \eqref{eq:plan1-dilL1} (after redefining $C_{\rm dil}$ to
absorb the factor of $2$). Estimate \eqref{eq:plan1-dilAlpha} follows
from \eqref{eq:bdv42ii}; note $\lambda U_*\subset B_{3R/2}$ so we apply
\eqref{eq:bdv42ii} on $B_{3R/2}$, picking up $C_2(3R/2)\le C_2(R)$
after enlarging $C_2$ by a constant depending on $R$ (still
$=:C_2(R)$).
\end{proof}

We now state the central quotient-preservation lemma.

\begin{lemma}[Lemma~3: quotient preservation after volume normalisation]
\label{lem:plan1-sel-L3}
Set $\lambda:=r_*^{-1}=(|B_1|/|U_*|)^{1/n}$ and
$\Otilde:=\lambda U_*$. Under the smallness regime
\eqref{eq:plan1-Sml} and with $\tau=q^{1/4}$, after translation by
$-x_{U_*}$ if needed:
\begin{enumerate}
\item[\rm(a)] $|\Otilde|=|B_1|=\omega_n$;
\item[\rm(b)] $|\BDValpha(\Otilde)-\eps|\le \rho\,\eps$ where
\begin{equation}\label{eq:plan1-rho}
\rho=\rho(q)
:=\frac{\sqrt{2q+q^{2}}}{\tau}+C_{\BDValpha}(n,R)\,q
\;\le\;2q^{1/4}+C_{\BDValpha}(n,R)\,q;
\end{equation}
for $q\le q_{\rm sel}(n,R)$, $\rho(q)\le 1/2$;
\item[\rm(c)] $E(\Otilde)-E(B_1)\le C_E(n,R)\,\delta$;
\item[\rm(d)] $\displaystyle Q_{\BDValpha}(\Otilde)\;\le\;
\frac{C_E(n,R)}{1-\rho(q)}\,Q_{\BDValpha}(\Om_0)$.
\end{enumerate}
The constants $C_{\BDValpha}(n,R)$ and $C_E(n,R)$ are given
explicitly in \eqref{eq:plan1-Calpha-def} and
\eqref{eq:plan1-CE-def} below.
\end{lemma}

\begin{proof}
\emph{Part~(a).} By construction
$|\Otilde|=\lambda^{n}|U_*|=(\omega_n/|U_*|)|U_*|=\omega_n$.

\emph{Part~(b): asymmetry control.} Combining
\eqref{eq:plan1-dilAlpha} with $|\lambda-1|\le 2C_5(n,R)\,\delta$ from
Lemma~\ref{lem:plan1-sel-L2} (which uses
$(\omega_n/s)^{1/n}-1$ has Lipschitz constant
$\le 2(n\omega_n)^{-1}$ near $s=\omega_n$),
\[
|\BDValpha(\Otilde)-\BDValpha(U_*)|
\;\le\;C_2(R)\,C_{\rm dil}(n,R)\cdot 2C_5(n,R)\,\delta
\;=:\;C_{\BDValpha}(n,R)\,\delta,
\]
i.e.\
\begin{equation}\label{eq:plan1-Calpha-def}
C_{\BDValpha}(n,R):=2\,C_2(R)\,C_{\rm dil}(n,R)\,C_5(n,R).
\end{equation}
Combined with \eqref{eq:plan1-L1-asym} and $\delta=\eps q$,
\[
|\BDValpha(\Otilde)-\eps|
\;\le\;|\BDValpha(\Otilde)-\BDValpha(U_*)|
+|\BDValpha(U_*)-\eps|
\;\le\;\eps\,C_{\BDValpha}(n,R)\,q
+\frac{\sqrt{2\eps\delta+\delta^2}}{\tau}.
\]
The square-root term equals
$\eps\,\sqrt{2q+q^{2}}/\tau$. Hence $\rho(q)$ defined in
\eqref{eq:plan1-rho}. With $\tau=q^{1/4}$ and $q\le1$,
$\sqrt{2q+q^{2}}/\tau\le \sqrt{3q}/q^{1/4}\le2q^{1/4}$. Under
\eqref{eq:plan1-Sml},
$\rho(q)\le 2 q_{\rm sel}^{1/4}+C_{\BDValpha}(n,R)\,q_{\rm sel}$;
choosing $q_{\rm sel}(n,R)$ small enough (still depending only on
$(n,R)$) gives $\rho(q)\le 1/2$.

\emph{Part~(c): energy scaling.} For every quasi-open
$V\subset B_R$ and $\mu>0$ with $\mu V\subset B_{R'}$,
\begin{equation}\label{eq:plan1-Escale}
E(\mu V)=\mu^{n+2}E(V).
\end{equation}
\emph{Proof of \eqref{eq:plan1-Escale}.} If $u_V$ is the torsion
function of $V$, define $w(x):=\mu^{2}u_V(x/\mu)$ for $x\in\mu V$.
Then $\nabla w(x)=\mu\nabla u_V(x/\mu)$, $\Delta w(x)=\Delta u_V(x/\mu)
=-1$, $w=0$ on $\partial(\mu V)$, so $w=u_{\mu V}$. Substituting,
\[
E(\mu V)=-\tfrac12\!\int_{\mu V}u_{\mu V}\,dx
=-\tfrac12\!\int_{\mu V}\mu^{2}u_V(x/\mu)\,dx
=-\tfrac12\,\mu^{n+2}\!\int_V u_V(y)\,dy=\mu^{n+2}E(V).
\]
Apply \eqref{eq:plan1-Escale} with $\mu=\lambda=r_*^{-1}$:
\begin{equation}\label{eq:plan1-Etilde-decomp}
E(\Otilde)-E(B_1)
=\lambda^{n+2}\bigl(E(U_*)-E(B_1)\bigr)
+(\lambda^{n+2}-1)\,E(B_1).
\end{equation}
From \eqref{eq:plan1-L1-energy} and
\eqref{eq:plan1-Lem45},
\[
E(U_*)-E(B_1)
\;\le\;F_{\hat\eta}(U_*)-F_{\hat\eta}(B_1)+\hat\eta^{-1}||U_*|-\omega_n|
\;\le\;\bigl(1+\hat\eta^{-1}C_4(n,R)\bigr)\delta
\;=:\;C_6(n,R)\,\delta.
\]
By Lemma~\ref{lem:plan1-sel-L2} and $|\lambda-1|\le1/2$,
\[
\lambda^{n+2}\le (3/2)^{n+2}=:C_7(n),
\qquad
|\lambda^{n+2}-1|\le (n+2)(3/2)^{n+1}|\lambda-1|
\le (n+2)(3/2)^{n+1}C_5(n,R)\,\delta=:C_8(n,R)\,\delta.
\]
Since $|E(B_1)|=T_n/2=\omega_n/(2n(n+2))$ is dimensional, plugging
into \eqref{eq:plan1-Etilde-decomp},
\begin{align*}
E(\Otilde)-E(B_1)
&\le C_7(n)\,C_6(n,R)\,\delta + C_8(n,R)\,|E(B_1)|\,\delta\\
&= \bigl(C_6 C_7+C_8|E(B_1)|\bigr)\,\delta
\;=:\;C_E(n,R)\,\delta,
\end{align*}
with
\begin{equation}\label{eq:plan1-CE-def}
C_E(n,R):=
\bigl(1+\hat\eta(n,R)^{-1}C_4(n,R)\bigr)\,(3/2)^{n+2}
+\frac{(n+2)(3/2)^{n+1}\,C_5(n,R)\,\omega_n}{2n(n+2)}.
\end{equation}
This is \eqref{eq:plan1-CE-def}; we have used $E(B_1)<0$ when
applying \eqref{eq:plan1-Etilde-decomp}; the signed contribution
$(\lambda^{n+2}-1)E(B_1)$ has the wrong sign when $\lambda>1$ but is
controlled in absolute value.

\emph{Part~(d): quotient.} Combining (b) and (c) with
$\BDValpha(\Otilde)\ge(1-\rho(q))\eps$,
\[
Q_{\BDValpha}(\Otilde)
=\frac{E(\Otilde)-E(B_1)}{\BDValpha(\Otilde)}
\;\le\;\frac{C_E(n,R)\,\delta}{(1-\rho(q))\,\eps}
=\frac{C_E(n,R)}{1-\rho(q)}\,Q_{\BDValpha}(\Om_0).
\]
This completes the proof.
\end{proof}

\begin{remark}[Preservation of the deficit-to-asymmetry ratio]
\label{rmk:plan1-sel-ratio}
Lemma~\ref{lem:plan1-sel-L3}(d) is the load-bearing
deficit-to-asymmetry preservation statement: starting from
$Q_{\BDValpha}(\Om_0)=q$, the volume-normalised companion satisfies
$Q_{\BDValpha}(\Otilde)\le \widetilde C(n,R)\,q$ for an explicit constant
$\widetilde C(n,R)=2C_E(n,R)$ once $\rho(q)\le 1/2$. The improvement
factor on the asymmetry is
$\BDValpha(\Otilde)\ge \tfrac12\BDValpha(\Om_0)$, so the new asymmetry
is at most a factor of $2$ smaller than the old. The chain
\eqref{eq:plan1-Frbound} then transfers the same statement to the
Fraenkel quotient at the price of a square root.

Volume normalisation enters the proof in two distinct places:
\begin{itemize}
\item In \emph{(b)}, the volume defect produces the
$|\lambda-1|\le C_5\delta$ contribution to the asymmetry change
through Lemma~\ref{lem:plan1-sel-dilation}.
\item In \emph{(c)}, the volume-scaling identity
\eqref{eq:plan1-Escale} produces the explicit
$\lambda^{n+2}$-factor, which is the only place where the dimension
$n$ enters the constant $C_E(n,R)$.
\end{itemize}
Both contributions are $O(\delta)$, hence subleading compared to the
$O(\eps q^{1/4})$ contribution from
\eqref{eq:plan1-L1-asym}. This is the precise mechanism by which the
single-set selection step preserves the quotient $Q_{\BDValpha}$ up to a
multiplicative $(n,R)$-constant.
\end{remark}

\medskip

\subsubsection{The selection theorem and BDV regularity package}
\label{subsubsec:plan1-sel-thm}

\begin{theorem}[Single-set quantitative selection principle]
\label{thm:plan1-sel-main}
Fix $n\ge2$ and $R\ge2$. There exist constants
\[
\eps_{\rm sel}(n,R)>0,\quad q_{\rm sel}(n,R)>0,\quad
C_{\rm sel}(n,R)>0,\quad P_{\rm sel}(n,R)>0,
\]
depending only on $(n,R)$ via the named BDV constants of
\eqref{eq:plan1-tauex} and \eqref{eq:plan1-CE-def},
with the following property. Let $\Om_0\subset B_R$ be quasi-open
with
\[
|\Om_0|=\omega_n,\qquad
0<\eps:=\BDValpha(\Om_0)\le\eps_{\rm sel}(n,R),\qquad
0<q:=Q_{\BDValpha}(\Om_0)\le q_{\rm sel}(n,R).
\]
Set $\tau:=q^{1/4}$, let $U_*$ be a minimiser of the functional $J$ in
\eqref{eq:plan1-Jdef} provided by
Lemma~\ref{lem:plan1-sel-existence}, and define
\[
r_*:=(|U_*|/\omega_n)^{1/n},
\qquad
\Otilde:=r_*^{-1}U_*.
\]
Then, after the translation $\Otilde\mapsto \Otilde-x_{\Otilde}$ that
brings the barycentre to the origin:
\begin{enumerate}
\item[\rm(i)] $P(U_*)\le P_{\rm sel}(n,R)$ and $|\Otilde|=\omega_n$;
\item[\rm(ii)] $(1-\rho)\eps\le \BDValpha(\Otilde)\le (1+\rho)\eps$ with
$\rho=\rho(q)\le 1/2$ given by \eqref{eq:plan1-rho}, hence
$\tfrac12\,\BDValpha(\Om_0)\le \BDValpha(\Otilde)\le \tfrac32\,\BDValpha(\Om_0)$;
\item[\rm(iii)] $\displaystyle
Q_{\BDValpha}(\Otilde)\;\le\;
\frac{C_{\rm sel}(n,R)}{1-\rho(q)}\,Q_{\BDValpha}(\Om_0)
\;\le\; 2\,C_{\rm sel}(n,R)\,Q_{\BDValpha}(\Om_0)
$, with $C_{\rm sel}(n,R):=C_E(n,R)$ from \eqref{eq:plan1-CE-def};
\item[\rm(iv)] $\Otilde$ inherits the BDV penalised-minimiser regularity
package on its torsion function $\tilde u:=\tilde u_{\Otilde}$ via the
perturbed quasi-minimality inequality \eqref{eq:plan1-qmin} below: namely
\begin{itemize}
\item two-sided density estimates on $\Otilde$ at scales
$\le r_0=r_0(n,R)$ (BDV Lemma~4.9);
\item Lipschitz/nondegeneracy of $\tilde u$ near $\partial\Otilde$ (BDV
Lemma~4.11, 4.12);
\item smooth, uniformly $C^{k}(n,R)$-bounded Bernoulli coefficient on
$\partial\Otilde$ (BDV Lemma~4.16);
\end{itemize}
all constants depending only on $(n,R)$;
\item[\rm(v)] On $\partial\Otilde$ the selected Bernoulli law takes the
volume-rescaled form of Lemma~\ref{lem:plan1-sel-bern} below.
\end{enumerate}
\end{theorem}

\begin{proof}
Parts (i)--(iii) are
Lemmas~\ref{lem:plan1-sel-existence}--\ref{lem:plan1-sel-L3}.
Parts (iv)--(v) are stated and justified next.
\end{proof}

\paragraph{Perturbed quasi-minimality inequality.}
By a direct comparison test in $J$, the minimiser $U_*=\{u_*>0\}$ (where
$u_*=u_{U_*}$ is the torsion function) satisfies the BDV (4.19)
inequality: for every $v\in H^1_0(B_R)$,
\begin{equation}\label{eq:plan1-qmin}
\tfrac12\!\int|\nabla u_*|^2-\!\int u_*
+f_{\hat\eta}(|\{u_*>0\}|)
\;\le\;\tfrac12\!\int|\nabla v|^2-\!\int v+f_{\hat\eta}(|\{v>0\}|)
+C_R\,\tau\,\bigl|\{u_*>0\}\Delta\{v>0\}\bigr|.
\end{equation}
The error term arises from the Lipschitz bound on the square-root
penalty in \eqref{eq:plan1-Jdef}, dominated by
$\tau\,C_2(R)|U\Delta U_*|$, where $C_2(R)$ is the BDV
constant of \eqref{eq:bdv42ii}; so $C_R$ may be taken equal to
$C_2(R)$. For $\tau\le\tau_{\rm reg}(n,R)$,
\eqref{eq:plan1-qmin} places $u_*$ in the
Alt--Caffarelli quasi-minimiser class with smallness constant
controlled by $(n,R)$, and the BDV
package \cite[Lemmas~4.9, 4.11, 4.12, 4.16]{BDV} applies; this yields
the density, Lipschitz, nondegeneracy and Bernoulli statements of
Theorem~\ref{thm:plan1-sel-main}(iv). The transfer from $U_*$ to
$\Otilde=r_*^{-1}U_*$ is a bilipschitz dilation with
$|\lambda-1|\le 1/2$ (Lemma~\ref{lem:plan1-sel-L2}); the density,
Lipschitz, and $C^{k}$ constants degrade by a factor at most
$(3/2)^{n}$, which is absorbed into the $(n,R)$-bookkeeping.

\paragraph{Volume-rescaled Bernoulli law.}

\begin{lemma}[Volume-rescaled Bernoulli law on $\partial\Otilde$]
\label{lem:plan1-sel-bern}
Let $\tilde u(y):=r_*^{-2}u_{U_*}(r_*y)$ for $y\in\Otilde=r_*^{-1}U_*$.
Then $-\Delta\tilde u=1$ in $\Otilde$, $\tilde u=0$ on $\partial\Otilde$,
and on $\partial\Otilde$ the BDV Lemma~4.16 Bernoulli law takes the form
\begin{equation}\label{eq:plan1-Berntilde}
|\nabla\tilde u(y)|^2
=\tilde\Lambda
+\tilde a_\sigma\,\tilde\Psi_{r_*}(y),
\end{equation}
where
\begin{equation}\label{eq:plan1-Berncoeff}
\tilde\Lambda=r_*^{-2}\Lambda,
\qquad
\tilde a_\sigma=r_*^{-1}\,a_\sigma,
\qquad
|a_\sigma|\le \tau,
\end{equation}
and
\[
\tilde\Psi_{r_*}(y)
:=|y-r_*^{-1}x_{U_*}|-r_*^{n}\,V_{\Otilde}\cdot y,
\qquad
V_{\Otilde}:=\!\int_{\Otilde}\!\frac{z-x_{\Otilde}}{|z-x_{\Otilde}|}\,dz.
\]
Both prefactors satisfy
$\tilde\Lambda=\Lambda\,(1+O(\delta))$ and
$\tilde a_\sigma=a_\sigma\,(1+O(\delta))$ as $\delta\to0$, with
implicit constants depending only on $(n,R)$.
\end{lemma}

\begin{proof}
For $y\in\Otilde$, $\nabla\tilde u(y)=r_*^{-1}\nabla u_{U_*}(r_*y)$ and
$\Delta\tilde u(y)=\Delta u_{U_*}(r_*y)=-1$, so $\tilde u$ solves
$-\Delta\tilde u=1$ in $\Otilde$ with zero Dirichlet condition on
$\partial\Otilde$ (which is $r_*^{-1}\partial U_*$). Squaring,
$|\nabla\tilde u(y)|^2=r_*^{-2}|\nabla u_{U_*}(r_*y)|^2$.

By BDV Lemma~4.16, on $\partial U_*$,
$|\nabla u_{U_*}(x)|^2=\Lambda+a_\sigma\Psi_{U_*}(x)$ with
$\Psi_{U_*}(x)=|x-x_{U_*}|-V_{U_*}\cdot x$ and
$V_{U_*}=\int_{U_*}(z-x_{U_*})/|z-x_{U_*}|\,dz$. Plugging $x=r_*y$,
\[
|\nabla\tilde u(y)|^2
=r_*^{-2}\Lambda+r_*^{-2}a_\sigma\Psi_{U_*}(r_*y).
\]
Change of variables $z=r_*w$ in the centroid integral:
\[
V_{U_*}
=\!\int_{U_*}\frac{z-x_{U_*}}{|z-x_{U_*}|}dz
=r_*^{n}\!\int_{\Otilde}\frac{w-x_{\Otilde}}{|w-x_{\Otilde}|}dw
=r_*^{n}V_{\Otilde},
\]
using $x_{U_*}=r_*x_{\Otilde}$ for the barycentre. Hence
\[
\Psi_{U_*}(r_*y)
=|r_*y-x_{U_*}|-V_{U_*}\cdot r_*y
=r_*\bigl(|y-r_*^{-1}x_{U_*}|-r_*^{n}V_{\Otilde}\cdot y\bigr)
=r_*\,\tilde\Psi_{r_*}(y).
\]
Substituting,
\(
|\nabla\tilde u(y)|^2
=r_*^{-2}\Lambda+r_*^{-1}a_\sigma\,\tilde\Psi_{r_*}(y),
\)
which is \eqref{eq:plan1-Berntilde} with the coefficients
\eqref{eq:plan1-Berncoeff}. Lemma~\ref{lem:plan1-sel-L2} gives
$|r_*-1|\le C_5(n,R)\delta$, hence
$\tilde\Lambda=\Lambda(1+O_{n,R}(\delta))$ and analogously for
$\tilde a_\sigma$.
\end{proof}

\subsubsection{Summary of constants and dependences}
\label{subsubsec:plan1-sel-constants}

For ease of cross-reference in later subsections, we record the
constants used in Theorem~\ref{thm:plan1-sel-main} together with their
dependences. All constants depend only on the pair $(n,R)$ via named BDV
constants; none is set by a compactness or contradiction step.

\begin{center}
\renewcommand{\arraystretch}{1.20}
\begin{tabular}{|l|l|l|}
\hline
Symbol & Defined in / source & Depends on \\\hline
$\hat\eta,C_4$ & BDV Lemma~4.5 & $n,R$ \\
$C_1$ & BDV Lemma~4.2(i), eq.~\eqref{eq:bdv42i} & $n$ \\
$C_2$ & BDV Lemma~4.2(ii), eq.~\eqref{eq:bdv42ii} & $R$ \\
$\sigma_2,\sigma_3$ & BDV Lemmas~4.9, 4.16 & $n,R$ \\
$\tau_{\rm ex},\tau_{\rm reg}$ & eq.~\eqref{eq:plan1-tauex} & $n,R$ \\
$P_{\rm sel}$ & Lemma~\ref{lem:plan1-sel-existence} & $n,R$ \\
$C_5$ & Lemma~\ref{lem:plan1-sel-L2} & $n,R$ \\
$C_{\rm dil}$ & Lemma~\ref{lem:plan1-sel-dilation},
eq.~\eqref{eq:plan1-dilL1} & $n,R$ \\
$C_{\BDValpha}=2 C_2 C_{\rm dil}C_5$ & eq.~\eqref{eq:plan1-Calpha-def} & $n,R$ \\
$C_E$ & eq.~\eqref{eq:plan1-CE-def} & $n,R$ \\
$C_{\rm sel}=C_E$ & Theorem~\ref{thm:plan1-sel-main} & $n,R$ \\
$\eps_{\rm sel},q_{\rm sel}$ & eq.~\eqref{eq:plan1-Sml} & $n,R$ \\
$\rho(q)=\sqrt{2q+q^2}/\tau+C_{\BDValpha}q$ &
eq.~\eqref{eq:plan1-rho} & $n,R,q$ \\
$\tilde\Lambda=r_*^{-2}\Lambda$,
$\tilde a_\sigma=r_*^{-1}a_\sigma$ &
Lemma~\ref{lem:plan1-sel-bern} & $n,R$ \\
\hline
\end{tabular}
\end{center}

\paragraph{Outputs feeding the next subsection.}
The selection principle (Theorem~\ref{thm:plan1-sel-main}) produces a
volume-normalised quasi-open set $\Otilde\subset B_{R'(n,R)}$ with
$|\Otilde|=\omega_n$, $\BDValpha(\Otilde)\in[\eps/2,3\eps/2]$,
$Q_{\BDValpha}(\Otilde)\le 2C_{\rm sel}(n,R)\,Q_{\BDValpha}(\Om_0)$, and
the full BDV free-boundary regularity package. This is the input for
the graph-entry subsection (Plan 1 \S 3.2), which converts smallness of
$\BDValpha(\Otilde)$ into $C^{1,\gamma}$ normal-graph parametrisation
over $\partial B_1(x_{\Otilde})$ via the
$\alpha\to$Hausdorff$\to$flatness chain.
